\par
\Needspace{35\baselineskip}
\subsubsection*{Upright labels for phase and gain margins}
The template defines \verb|\PM| and \verb|\GM| in
\path{styles/math/config.tex}, loaded by
\path{styles/math/packages.tex}. They are template helpers,
not commands supplied by \texttt{amsmath}. Use them in math mode
to keep multi-letter labels upright rather than typesetting them
as products of italic variables.

\partlabel{Preamble -- definitions already provided by the template}
\begin{referencecode}
\newcommand{\PM}{\mathrm{PM}}
\newcommand{\GM}{\mathrm{GM}}
\end{referencecode}
When using the template, do not repeat these definitions in
\texttt{main.tex}. In a standalone document, add them to the preamble.

\partlabel{Margins, subscripts and an infinite gain margin}
\begin{referencecode}
\[
    \PM = 180^\circ + \angle L(j\omega_{gc})
    \qquad \GM = \frac{1}{|L(j\omega_{pc})|}
\]
\[
    \GM_{\mathrm{dB}} = 20\log_{10}(\GM)
    \qquad \GM = \infty
\]
\end{referencecode}

\[
    \PM = 180^\circ + \angle L(j\omega_{gc})
    \qquad \GM = \frac{1}{|L(j\omega_{pc})|}
\]
\[
    \GM_{\mathrm{dB}} = 20\log_{10}(\GM)
    \qquad \GM = \infty
\]
Here \(\omega_{gc}\) and \(\omega_{pc}\) denote gain- and
phase-crossover frequencies; the infinite margin illustrates the
case with no phase crossover. The gain margin \(\GM\) is a ratio;
\(\GM_{\mathrm{dB}}\) is its value in decibels. For numerical
angles and units, see \texttt{siunitx} on page~\pageref{units}.
\par
